\begin{prob}
    In each of the problems below state the best case and worst case time
    complexities of the given piece of code using asymptotic notation. Note that
    some algorithms may have the same best case and worst case time complexities. If
    the best and worst case complexities \emph{are} different, identify which
    inputs result in the best case and worst case. You do not otherwise need to show your work for
    this problem.

    \emph{Example Algorithm:} \python{linear_search} as given in lecture. 

    \emph{Example Solution:}
    Best case: $\Theta(1)$, when the target is the first element of the array.
    Worst case: $\Theta(n)$, when the target is not in the array.

    \begin{subprobset}
        \begin{subprob}
            \begin{minted}[autogobble]{python}
                def has_duplicates(data):
                    """Checks for duplicates in `data`. `data` is an array of size n."""
                    n = len(data)
                    for i in range(n):
                        for j in range(i + 1, n):
                            if data[i] == data[j]:
                                return True
                    return False
            \end{minted}

            \begin{soln}
                Best case: $\Theta(1)$, when the first and second elements of the array
                are duplicates.

                Worst case: $\Theta(n^2)$, when there are no duplicates.

                You have to wonder... isn't there a faster approach?
            \end{soln}
        \end{subprob}

        \begin{subprob}
            \begin{minted}[autogobble]{python}
                def f_2(data):
                    """`data` is an array of n numbers"""
                    m = data[0]
                    c = 1
                    for ix in range(1, len(data)):
                        if c == 0:
                            m = data[ix]
                            c = 1
                        elif m == data[ix]:
                            c += 1
                            if c > len(data) // 2:
                                return m
                        else:
                            c -= 1

                    return m
            \end{minted}
            

            \begin{soln}
                Best case and worst case are both $\Theta(n)$. In the ``best case'',
                the loop exits early when \python{c > len(data) // 2}. But $c$ grows by
                at most 1 on each iteration, so we'd need to go through at least
                $n / 2$ iterations to exit, which is still $\Theta(n)$.

                By the way, this algorithm is called the Boyer-Moore majority algorithm,
                and can be used to efficiently find a majority element of a list (as long
                as there is one). Though note that you can reason about the
                algorithm's time complexity without actually knowing what it
                computes.
            \end{soln}
        \end{subprob}

        \begin{subprob}
            \begin{minted}[autogobble]{python}
                def median(numbers):
                    """computes the median. `numbers` is an array of n numbers"""
                    n = len(numbers)
                    for x in numbers:
                        less = 0
                        more = 0
                        for y in numbers:
                            if y <= x:
                                less += 1
                            if y >= x:
                                more += 1
                        if less >= n/2 and more >= n/2:
                            return x
            \end{minted}
            
            \begin{soln}
                Best case: $\Theta(n)$, when the first element of the array is a median.

                Worst case: $\Theta(n^2)$, when the first $n-1$ elements are not medians, and the last element of the array is the median.
            \end{soln}
        \end{subprob}

        \begin{subprob}
            \begin{minted}[autogobble]{python}
                def index_of_median(numbers):
                    """`numbers` is an array of size n"""
                    # the median() from above
                    m = median(numbers)
                    # the linear_search() from lecture
                    return linear_search(numbers, m)
            \end{minted}
            

            \begin{soln}
                Best case: $\Theta(n)$. This occurs when the median is the first element of the array,
                as then \python{median} takes $\Theta(n)$ and \python{linear_search} takes $\Theta(1)$,
                for a total of $\Theta(n)$.

                Worst case: $\Theta(n^2)$. This occurs when the median is the last element of the
                array. In this situation, \python{median} takes $\Theta(n^2)$ time and \python{linear_search}
                takes $\Theta(n)$ time, for a total of $\Theta(n^2)$ time.
            \end{soln}
        \end{subprob}

    \end{subprobset}


\end{prob}
